\section{One-Dimensional Heat Diffusion Problem}

The physical problem considered in this work is the one-dimensional heat diffusion equation. This equation describes how temperature evolves in time within a medium due to thermal conduction. Despite its apparent simplicity, it captures many of the fundamental mechanisms encountered in diffusion-dominated physical systems and constitutes one of the most widely used benchmark problems for Physics-Informed Neural Networks (PINNs).

Studying this problem in detail provides a controlled mathematical framework where the behavior of PINNs can be analyzed rigorously before extending the methodology to more complex nonlinear or multiphysics systems.

% =========================================================
\subsection{Governing Equation}
% =========================================================

The one-dimensional heat equation is written as:

\begin{equation}
\frac{\partial T}{\partial t}
=
\alpha
\frac{\partial^2 T}{\partial x^2}
\end{equation}

where:
\begin{itemize}
    \item $T(x,t)$ is the temperature field,
    \item $x$ is the spatial coordinate,
    \item $t$ is time,
    \item $\alpha$ is the thermal diffusivity coefficient.
\end{itemize}

The thermal diffusivity $\alpha$ characterizes the rate at which heat propagates through the material. Physically, large values of $\alpha$ correspond to faster temperature diffusion.

The spatial domain considered here is:

\begin{equation}
x \in [0,L]
\end{equation}

with $L$ the length of the domain.

% =========================================================
\subsection{Initial and Boundary Conditions}
% =========================================================

To obtain a well-posed problem, the governing equation must be supplemented with initial and boundary conditions.

The boundary conditions are homogeneous Dirichlet conditions:

\begin{equation}
T(0,t)=0
\end{equation}

\begin{equation}
T(L,t)=0
\end{equation}

which impose zero temperature at both ends of the domain.

The initial condition is chosen as:

\begin{equation}
T(x,0)=\sin\left(\frac{\pi x}{L}\right)
\end{equation}

This initial condition is particularly convenient because it is compatible with the boundary conditions and corresponds to a single spatial eigenmode of the diffusion operator.

% =========================================================
\subsection{Derivation of the Exact Solution}
% =========================================================

The exact analytical solution can be derived using the method of separation of variables.

We assume that the solution can be written as:

\begin{equation}
T(x,t)=X(x)\Theta(t)
\end{equation}

where:
\begin{itemize}
    \item $X(x)$ contains the spatial dependence,
    \item $\Theta(t)$ contains the temporal dependence.
\end{itemize}

Substituting this expression into the heat equation gives:

\begin{equation}
X(x)\frac{d\Theta}{dt}
=
\alpha
\Theta(t)\frac{d^2X}{dx^2}
\end{equation}

Dividing both sides by $X(x)\Theta(t)$ yields:

\begin{equation}
\frac{1}{\Theta}\frac{d\Theta}{dt}
=
\alpha
\frac{1}{X}\frac{d^2X}{dx^2}
\end{equation}

The left-hand side depends only on time while the right-hand side depends only on space. Therefore, both sides must be equal to the same constant, denoted $-\lambda$:

\begin{equation}
\frac{1}{\Theta}\frac{d\Theta}{dt}
=
-\lambda
\end{equation}

\begin{equation}
\alpha
\frac{1}{X}\frac{d^2X}{dx^2}
=
-\lambda
\end{equation}

This leads to two ordinary differential equations.

% =========================================================
\subsubsection{Temporal Equation}
% =========================================================

\begin{equation}
\frac{d\Theta}{dt}
=
-\lambda \Theta
\end{equation}

whose solution is:

\begin{equation}
\Theta(t)=Ae^{-\lambda t}
\end{equation}

% =========================================================
\subsubsection{Spatial Equation}
% =========================================================

\begin{equation}
\frac{d^2X}{dx^2}
+
\frac{\lambda}{\alpha}X
=
0
\end{equation}

Defining:

\begin{equation}
k^2=\frac{\lambda}{\alpha}
\end{equation}

the equation becomes:

\begin{equation}
\frac{d^2X}{dx^2}+k^2X=0
\end{equation}

whose general solution is:

\begin{equation}
X(x)=B\sin(kx)+C\cos(kx)
\end{equation}

Applying the boundary condition at $x=0$:

\begin{equation}
X(0)=0
\end{equation}

gives:

\begin{equation}
C=0
\end{equation}

Thus:

\begin{equation}
X(x)=B\sin(kx)
\end{equation}

Applying the second boundary condition at $x=L$:

\begin{equation}
X(L)=0
\end{equation}

leads to:

\begin{equation}
\sin(kL)=0
\end{equation}

which implies:

\begin{equation}
k_n=\frac{n\pi}{L}
\end{equation}

for integer values $n=1,2,3,\dots$

The associated eigenvalues are:

\begin{equation}
\lambda_n
=
\alpha
\left(\frac{n\pi}{L}\right)^2
\end{equation}

The general solution of the heat equation is therefore:

\begin{equation}
T(x,t)
=
\sum_{n=1}^{\infty}
A_n
\sin\left(\frac{n\pi x}{L}\right)
\exp\left[
-\alpha
\left(
\frac{n\pi}{L}
\right)^2
t
\right]
\end{equation}

% =========================================================
\subsection{Particular Solution Used in This Work}
% =========================================================

In this study, the initial condition corresponds exactly to the first eigenmode:

\begin{equation}
T(x,0)
=
\sin\left(\frac{\pi x}{L}\right)
\end{equation}

Therefore, only the first term of the series is present, yielding the exact solution:

\begin{equation}
T(x,t)
=
\sin\left(\frac{\pi x}{L}\right)
\exp\left[
-\alpha
\left(
\frac{\pi}{L}
\right)^2
t
\right]
\end{equation}

This analytical expression is particularly useful because:
\begin{itemize}
    \item it satisfies the governing equation exactly,
    \item it satisfies the boundary and initial conditions,
    \item it allows direct generation of synthetic datasets,
    \item and it provides an exact reference solution for evaluating PINN predictions and inverse parameter estimation performance.
\end{itemize}